\subsection{An integrated sequence of five experimental phases}
The study began with a specific intervention in a lightweight forecasting model: changing its decomposition and distinguishing an observed measurement from a carried or fallback value. This starting point is useful because an imputed scalar alone does not identify the measurement process that generated it. Nevertheless, representing that process increases capacity and can encourage a model to fit the synthetic generator. A successful missingness descriptor must therefore justify itself through both held-out error and computational measurements. The historical E0--E5 experiments established the initial empirical question; compact variants tested whether a smaller summary could preserve the apparent benefit; additional comparators and service cohorts subsequently tested its portability.

The five phases are complementary rather than independent replications. Phase I examines temporal decomposition and missingness heads on the fixed 28-service source cohort. Phase II evaluates C0--C3 and the validation-selected alias V2 on that same known test. Phase III refits source models under the expanded protocol and transfers them to 32 additional services. Phase IV stresses these models with additional mechanisms, direct horizons, longer histories and a separate short Alibaba case. Phase V measures computation in the actual implementations. The narrative follows the progression of the scientific question, while the chronology of test exposure remains part of the validity analysis.

In particular, the source-test observation that E5 helped memory motivated subsequent model development. The new-service protocol froze four primary contrasts before evaluating those destination outcomes, but the source period was already familiar. Consequently, the locally frozen contrasts improve procedural transparency without making the entire project prospectively blinded. An effect can be consistent across initializations of a fixed model and still fail to extend across services or datasets. Maintaining these distinctions is necessary to explain why compactness, seed consistency and transfer accuracy need not agree.

The research gap addressed here is a joint empirical accounting of decomposition, measurement provenance, capacity and transfer under one recorded pipeline. It is narrower than proposing a general solution to cloud resource management. Forecasting measured usage supplies one input to a future management decision, whereas a controller also requires a policy, capacity constraints and an action-dependent system response. The implemented experiments assess point error and execution cost. They cannot identify how unobserved demand would respond to allocating more resources, and they do not estimate the monetary value of a small MASE change.
